\documentclass[a4paper, 12pt]{article}
\usepackage[utf8]{inputenc}
\usepackage[T2A]{fontenc}
\usepackage[russian]{babel}
\usepackage{geometry}
\geometry{left=20mm,right=20mm,top=25mm,bottom=25mm}
\usepackage{enumitem}
\usepackage{xcolor}
\usepackage{amsmath}
\usepackage{amssymb}
\usepackage{tcolorbox}
\usepackage{titlesec}

\definecolor{questioncolor}{RGB}{0,60,120}
\definecolor{answercolor}{RGB}{20,90,20}
\definecolor{tipcolor}{RGB}{150,0,0}

\newcommand{\question}[1]{%
  \vspace{1.2em}
  \noindent\rule{\textwidth}{0.8pt}
  \vspace{0.3em}

  \noindent{\large\textbf{\textcolor{questioncolor}{Вопрос:}} #1}
  \vspace{0.5em}
}

\newcommand{\answer}[1]{%
  \noindent\textbf{\textcolor{answercolor}{Как отвечать:}} #1
}

\newcommand{\tip}[1]{%
  \vspace{0.3em}
  \noindent\textit{\textcolor{tipcolor}{Совет: #1}}
}

\title{\textbf{Вопросы комиссии на защите диплома}\\[0.5em]
{\large Vehicle Re-ID: ResNet-50 / ResNet-101 vs ViT-Small}\\[0.3em]
{\normalsize С подробными ответами и советами}}
\author{}
\date{}

\begin{document}
\maketitle
\tableofcontents
\newpage

% ================================================================
\section{Архитектурные вопросы}
% ================================================================

\question{1. У ResNet-50 параметров больше (25M), чем у ViT-Small (22M). Почему же ResNet требует больше эпох (120 против 50), но итоговое качество ниже?}

\answer{«Из-за отсутствия inductive bias трансформерам нужно больше времени и аккуратный warmup, чтобы выучить базовые паттерны --- например, трансляционную инвариантность, которая в CNN заложена архитектурно через свёртки. Но как только ViT обучается, его механизм глобального внимания позволяет связывать удалённые участки изображения (фары и форму крыши, логотип и диски) гораздо эффективнее, чем CNN через глубокую иерархию локальных рецептивных полей. Кроме того, ViT-Small предобучен на ImageNet-21k (14M изображений) против ImageNet-1k (1.2M) у ResNet-50 --- стартовые веса качественно лучше, поэтому fine-tuning быстрее.»}

\tip{Это стандартный вопрос. Упомяни ImageNet-21k pretrain --- это сильный аргумент.}

% ================================================================

\question{2. ResNet-101 показал результаты хуже ResNet-50 (Rank-1: 86.7\% против 88.7\%, mAP: 57.7\% против 64.3\%). Как объяснить деградацию более глубокой сети?}

\answer{«Проблема в соотношении объёма данных и числа параметров. VeRi-776 содержит 37\,746 обучающих изображений и 575 ID --- это мало для эффективного обучения 43.7M параметров ResNet-101. Для сравнения, ImageNet-1k содержит 1.2M изображений для 25M параметров --- соотношение в 10 раз лучше. Второй фактор: при одинаковых гиперпараметрах (одинаковый lr, одинаковые 120 эпох) более глубокая сеть получает меньше обновлений на параметр. Косвенным свидетельством является больший прирост от re-ranking у ResNet-101 (+2.4\% mAP против +3.2\% у ResNet-50) --- это признак менее структурированного пространства эмбеддингов.»}

\tip{Подготовься к уточняющему вопросу: «А как исправить?» --- ответ: уменьшить lr, добавить аугментации, заморозить нижние слои, использовать более сильный pretrain.}

% ================================================================

\question{3. Почему вы использовали ViT-Small, а не ViT-Base? Ведь TransReID (97.1\%) использует ViT-Base.}

\answer{«Ограничения GPU: GTX 1660 Super имеет только 6 GB VRAM. ViT-Base (86M параметров) при том же batch size просто не поместится. ViT-Small (22M) --- оптимальный выбор для данного железа. При переходе на ViT-Base исторически наблюдается прирост 5--7\% по mAP (TransReID, 2021). Наша работа не ставила целью побить SOTA --- цель: честное сравнение CNN и Transformer в одинаковых условиях на одном GPU.»}

% ================================================================

\question{4. Что такое CLS-токен и почему именно его используют как глобальный дескриптор, а не усреднение всех патчей?}

\answer{«CLS-токен (classification token) --- специальный обучаемый вектор, добавляемый к последовательности патчей перед подачей в Transformer. Через механизм self-attention он взаимодействует со всеми патчами на каждом слое и агрегирует информацию со всего изображения. В отличие от простого усреднения (average pooling), CLS-токен обучается \emph{как} агрегировать --- он учится придавать больший вес дискриминативным патчам. Это подтверждается нашими картами внимания: при верных совпадениях CLS-токен акцентирует внимание на номерных знаках и дисках, а не равномерно по всему изображению.»}

% ================================================================

\question{5. В чём разница между GradCAM (для ResNet) и Attention Maps (для ViT)? Можно ли сравнить эти визуализации напрямую?}

\answer{«GradCAM использует градиенты от целевого класса через feature map последнего свёрточного слоя. Это post-hoc объяснение, косвенно отражающее важность признаков. Attention Maps ViT --- это явные веса внимания из архитектуры, показывающие, на какие патчи смотрит CLS-токен при \emph{каждом} прямом проходе. Принципиальное отличие: ViT-карты доступны без backpropagation, в реальном времени. Напрямую сравнивать их сложно --- разные пространственные разрешения (feature map у ResNet $\sim$16$\times$8, патчи у ViT 14$\times$14) и разная семантика. В нашей работе мы показываем только attention maps для ViT и качественное описание локальных признаков для ResNet.»}

% ================================================================
\section{Вопросы про метрики и оценку}
% ================================================================

\question{6. Rank-1 у ViT-Small = 90.05\%, а mAP = 69.41\%. Почему такой большой разрыв? Разве они не должны быть близки?}

\answer{«Rank-1 измеряет только первое место --- нашли правильную машину первой? mAP учитывает все позиции всех правильных ответов. В VeRi-776 одна машина присутствует в галерее несколько раз (с разных камер). Если система нашла одно правильное изображение на 1-й позиции, но остальные 4 копии спрятала на 50--100-е места --- Rank-1 = 100\%, а mAP будет гораздо ниже. Разрыв в 20+ процентных пунктов типичен для задачи Vehicle ReID на VeRi-776 и отражает именно это явление: система хорошо находит \emph{одного} кандидата, но плохо ранжирует \emph{всех}.»}

% ================================================================

\question{7. Коэффициент силуэта у ViT-Small = 0.260, у ResNet-50 = 0.024. Это большая разница. Но при этом разница в Rank-1 всего 1.4\%. Как объяснить такое расхождение?}

\answer{«Коэффициент силуэта измеряет геометрию всего пространства эмбеддингов --- насколько хорошо разделены 200 классов в целом. Rank-1 измеряет локальный результат: попал ли правильный ответ на первое место при конкретном запросе. Даже при плохой глобальной структуре (ResNet, силуэт 0.024) ближайший сосед может оказаться правильным в большинстве случаев --- система "случайно угадывает" правильно для простых запросов. Разрыв проявляется на сложных случаях: темные машины, схожие формы --- там ViT выигрывает явно. Кроме того, mAP более чувствителен к структуре пространства: разница mAP 5.1 процентных пункта гораздо значительнее, чем 1.4 пункта Rank-1.»}

% ================================================================

\question{8. Протокол no-same-camera --- что это и зачем? Можно ли без него?}

\answer{«No-same-camera запрещает сопоставлять изображения одного и того же транспортного средства с одной камеры. Без этого протокола система могла бы "читерить": если запрос взят с камеры 5, а в галерее есть изображение той же машины с той же камеры 5 --- они будут практически идентичны (одинаковый ракурс, освещение). Метрики были бы завышены и не отражали бы реальную трудность задачи. Протокол имитирует реальный сценарий: мы видели машину на одной камере и ищем её на \emph{других} камерах. Именно поэтому у камеры C18 Rank-1 = 0\% --- там условия съёмки кардинально отличаются от всех остальных камер.»}

% ================================================================
\section{Вопросы про визуализацию и анализ}
% ================================================================

\question{9. Вы показываете карты внимания. Но как вы можете быть уверены, что модель действительно понимает, что такое «фара» или «диск», а не просто реагирует на контрастные пиксели?}

\answer{«ViT обучается без разметки частей автомобиля (без part-based аннотаций). Модель ищет наиболее дискриминативные паттерны. Однако, как видно на успешных примерах, эти паттерны физически совпадают с уникальными деталями: логотипы, решётки, наклейки, рисунок дисков. Это не случайно --- эти объекты действительно наиболее информативны для различения машин. Строго доказать понимание невозможно --- можно лишь показать корреляцию. Для строгого доказательства применяют метод окклюзии: закрываем hot-зону черным квадратом и наблюдаем падение уверенности. В нашей работе этот эксперимент не проводился --- это направление будущих исследований.»}

\tip{Честность про ограничения --- признак зрелости исследователя. Комиссия это ценит.}

% ================================================================

\question{10. Камера C18 показывает Rank-1 = 0\% для всех пар. Это баг или фича? Как вы проверили, что это не ошибка в коде?}

\answer{«Это особенность данных, а не баг кода. Несколько аргументов: во-первых, обе модели (ResNet-50 и ViT-Small) независимо показывают нулевое качество для C18 --- маловероятно, что баг затронул обе модели одинаково. Во-вторых, граф переходов показывает, что C18 имеет мало входящих рёбер --- она редко связана с другими камерами. Это согласуется с гипотезой о нетипичных условиях съёмки (другой ракурс, освещение, высота монтажа). В реальной системе видеонаблюдения такой результат --- сигнал к проверке и техническому обслуживанию этой камеры.»}

% ================================================================

\question{11. Граф переходов содержит 108 рёбер из 380 возможных. Откуда такая разреженность? Это нормально?}

\answer{«Абсолютно нормально и физически интерпретируемо. VeRi-776 снимался на реальной дорожной сети Тяньцзиня (Китай). Дорожная сеть --- это граф, и машина может ехать только по существующим дорогам. Если камеры 3 и 17 расположены в разных районах города без прямого соединения --- переход C3$\to$C17 невозможен физически. 108 рёбер из 380 (28\%) отражают топологию реальной дорожной сети. Именно это делает граф переходов полезным инструментом: он показывает структуру задачи, а не просто статистику.»}

% ================================================================
\section{Вопросы про методы и алгоритмы}
% ================================================================

\question{12. Почему расстояние Махаланобиса дало деградацию –31\% mAP для ViT, хотя теоретически оно <<лучше>> учитывает структуру данных?}

\answer{«L2-нормализация проецирует все векторы на поверхность единичной многомерной сферы $\mathcal{S}^{383}$. Косинусное расстояние измеряет угол между ними --- это геометрически корректная метрика \emph{на сфере}. Матрица Махаланобиса $\Sigma^{-1}$ описывает эллипсоид в евклидовом пространстве. Если применить эллиптическое искажение к точкам, жёстко лежащим на идеальной сфере --- мы ломаем их угловые отношения. Результат: ранжирование рассыпается. Для ResNet-50 использовалось диагональное приближение ($D=2048$ --- полная матрица требует $2048^2 = 4M$ значений и нестабильна) --- это эквивалентно нормировке на стандартное отклонение, эффект нулевой (+0.28\%). Вывод: для нормированных гиперсферических эмбеддингов косинусное расстояние оптимально.»}

\tip{За этот ответ с объяснением про гиперсферу --- почти гарантированная «пятёрка» от аналитика в комиссии.}

% ================================================================

\question{13. Вы применили re-ranking с параметрами $k_1=20$, $k_2=6$, $\lambda=0.3$. Как вы выбрали эти значения? Проводили ли grid search?}

\answer{«Параметры взяты из оригинальной статьи Zhong et al. (2017) --- авторы k-reciprocal re-ranking показали, что эти значения оптимальны для задач ReID на стандартных бенчмарках. В нашей работе grid search не проводился из-за ограничений по времени вычислений --- re-ranking $O(N^2)$ дорог при $N=11\,579$ изображений галереи. Следующий шаг --- систематический поиск по сетке на validation split, что потенциально может дать ещё 1--2\% прироста mAP.»}

% ================================================================

\question{14. Почему метод Attention Rollout дал плохие карты внимания, и что вы использовали вместо него?}

\answer{«Attention Rollout перемножает усреднённые по головам матрицы внимания всех $L=12$ блоков. При умножении большого числа сглаженных матриц (каждая из которых уже частично равномерна из-за softmax) произведение быстро сходится к равномерному распределению --- дискриминативный сигнал теряется. Вместо этого мы используем агрегацию внимания только из последних $K=4$ блоков с операцией максимума по головам. Операция максимума сохраняет сигнал наиболее «сфокусированных» голов и не размывает его усреднением. Дополнительно применяется перцентильная нормализация (5-й и 95-й перцентили) для предотвращения <<выжигания>> карты единичными выбросами.»}

% ================================================================
\section{Продуктовые и практические вопросы}
% ================================================================

\question{15. Вы показали прирост от k-reciprocal re-ranking. А вы можете внедрить эту систему на перекрёстке для работы в реальном времени?}

\answer{«Пайплайн делится на две части. Быстрая: YOLOv11n детектирует ТС ($\approx$100 fps), ViT-Small извлекает эмбеддинг (миллисекунды на батч) --- эту часть можно ставить в real-time. Медленная: k-reciprocal re-ranking пересчитывает попарные дистанции всей галереи --- операция $O(N^2)$, для 11\,579 изображений это дорого. Поэтому архитектура двухрежимная: в реальном времени используется FAISS с косинусным расстоянием, re-ranking применяется только для offline-криминалистики --- когда нужно максимально точно найти машину по архиву за сутки.»}

% ================================================================

\question{16. Что произойдёт с системой, если появится машина, которой не было в обучающей выборке (open-set scenario)?}

\answer{«ReID --- принципиально open-set задача. Во время тестирования все 200 тестовых ID не пересекаются с 575 обучающими ID. Поэтому система изначально обучена находить схожие, а не «узнавать» конкретные машины. Проблемы возникают, если новая машина слишком непохожа на все обучающие (экзотический цвет, редкий тип кузова) --- тогда эмбеддинг будет размещён плохо в пространстве. Решение: периодическое дообучение на новых данных или unsupervised domain adaptation на целевом распределении.»}

% ================================================================

\question{17. Как вы оцениваете скорость инференса? Подходит ли ViT-Small для реального времени?}

\answer{«На GPU (NVIDIA T4 / GTX 1660 Super) инференс ViT-Small: $\approx$2--5 мс на одно изображение при batch=1, $\approx$0.5 мс на изображение при batch=32. YOLOv11n добавляет ещё $\approx$5--10 мс на кадр. Итого пайплайн: $\approx$10--15 мс на кадр --- это 65--100 fps, что покрывает стандартные требования к видеонаблюдению (25--30 fps). Поиск по галерее с FAISS: $\approx$1 мс для 10\,000 векторов. Система реального времени технически осуществима.»}

% ================================================================
\section{Философские и концептуальные вопросы}
% ================================================================

\question{18. Если я закрою номерной знак на изображении машины --- насколько упадёт точность системы? Вы это проверяли?}

\answer{«Напрямую этот эксперимент не проводился --- это направление будущих исследований. Косвенно ответ даёт анализ attention maps: при верных совпадениях CLS-токен часто акцентирован на номерном знаке. Если его закрыть --- часть дискриминативного сигнала потеряется. Насколько упадёт точность --- зависит от того, насколько машина отличима по другим признакам (цвет, форма, диски). Для «обезличенной» реидентификации (без номеров) это важный вопрос приватности, актуальный для реального внедрения.»}

% ================================================================

\question{19. В чём принципиальное отличие вашей задачи от задачи верификации лиц (face verification)?}

\answer{«Общее: обе задачи --- метрическое обучение, поиск по эмбеддингам, протокол no-same-camera/no-same-session. Различия:
(1) Машины не имеют уникальных биометрических признаков --- в отличие от лица, автомобиль серийный. Два одинаковых автомобиля одной модели/цвета различить крайне сложно.
(2) Машины менее инвариантны к ракурсу: спереди и сзади выглядят принципиально по-разному (разные фары, бампер, багажник).
(3) Нет эквивалента 3D-модели лица --- для автомобилей часть информации (форма крыши) стабильна, часть (фары) зависит от ракурса.
(4) Задача масштабируется плохо: миллионы машин в городе против тысяч лиц в базе доступа.»}

% ================================================================

\question{20. Как бы вы расширили работу, если бы получили доступ к живым камерам Санкт-Петербурга?}

\answer{«Несколько уровней:
(1) \emph{Адаптация}: дообучить/протестировать модель на питерских данных --- явная проверка domain shift из VeRi (Китай) в Россию (разные марки авто, разные дорожные условия).
(2) \emph{Граф камер}: построить реальный граф переходов с GPS-координатами --- можно добавить временну\'{ю} информацию (машина не может пересечь город за 1 секунду).
(3) \emph{Трекинг}: добавить multi-object tracking (ByteTrack, StrongSORT) внутри одной камеры для построения траекторий.
(4) \emph{Re-ID через время}: искать машину не только в текущем кадре, но в архиве за сутки.
Это выходит за рамки дипломной работы --- задача для полноценного R\&D проекта.»}

% ================================================================

\question{21. Вы упоминаете DINOv2 и CLIP в разделе будущих направлений. В чём их принципиальное преимущество перед вашим подходом?}

\answer{«DINOv2 (Meta, 2023) --- self-supervised претрейн ViT на 142M изображениях без разметки. Признаки DINOv2 обладают явной семантической структурой (attention maps корреспондируют с семантическими частями объекта) без fine-tuning. Для ReID это означает: даже без обучения на VeRi эмбеддинги DINOv2 группируют похожие машины. CLIP (OpenAI) добавляет текстовое выравнивание: можно искать машину по описанию («красный внедорожник с царапиной на двери»). CLIP-ReID (2023) достигает 97.5\% Rank-1 именно за счёт этого. Наш ViT-Small pretrain на ImageNet-21k --- значительно слабее, но воспроизводим на студенческом GPU.»}

% ================================================================

\question{22. Вы обучили модель на VeRi-776. Если её использовать на американских или европейских дорогах --- как изменится качество?}

\answer{«VeRi-776 снят в Тяньцзине (Китай, 2016). Доменный сдвиг будет значительным:
(1) \emph{Марки автомобилей}: в Китае популярны Geely, BYD, Great Wall --- редкие в Европе. Эти ID не представлены в обучающей выборке.
(2) \emph{Дорожная инфраструктура}: разметка дорог, освещение, угол камер отличаются.
(3) \emph{Погодные условия}: VeRi снят преимущественно в хорошую погоду, без снега/тумана.
Ожидаемое падение качества: 10--20\% Rank-1. Решение: UDA (Unsupervised Domain Adaptation) или дообучение на нескольких изображениях из целевого домена (few-shot adaptation).»}

% ================================================================

\question{23. Что показывает анализ intra-class дисперсии и зачем он важен?}

\answer{«Intra-class дисперсия --- среднее L2-расстояние от эмбеддингов одного транспортного средства до центроида его кластера. Показывает: насколько сильно разброс условий съёмки (разные камеры, разные ракурсы) влияет на представление одной и той же машины. Результаты: у ViT-Small внутриклассовая дисперсия в 1.76 раза меньше (0.247 против 0.434 у ResNet-50). Это означает: ViT-Small более инвариантен к изменению условий съёмки --- одна и та же машина с 20 разных камер представлена более согласованно. Именно поэтому re-ranking даёт ViT-Small больший прирост: его ближайшие соседи более надёжны.»}

% ================================================================
\section{Самый сложный вопрос}
% ================================================================

\question{24. Ваш метод извлечения attention maps использует максимум по головам, а не среднее. Не теряете ли вы информацию о том, что делают остальные головы?}

\answer{«Да, намеренно. В трансформерах разные головы специализируются на разных типах зависимостей: одни --- на локальных паттернах, другие --- на глобальных структурах, третьи --- на цвете, четвёртые --- на форме. Среднее по всем головам смешивает эти сигналы в равных пропорциях и часто даёт равномерно <<тёплую>> карту --- не выделяющую ничего. Максимум по головам сохраняет сигнал наиболее фокусированной головы для каждого патча --- ту голову, которая максимально уверена в важности данного патча. Это heuristic без теоретических гарантий, но практически он показывает более острые и интерпретируемые карты. Альтернативы: Attention Rollout (мы объясняем почему отказались), DINO-style (специальный регуляризатор на этапе обучения).»}

\tip{Это вопрос от специалиста по трансформерам. Если такой человек в комиссии --- ответ выше покажет глубокое понимание.}

% ================================================================
\vspace{2em}
\noindent\rule{\textwidth}{1pt}
\begin{center}
\textit{Подготовлено для защиты дипломной работы по Vehicle Re-Identification}\\
\textit{ResNet-50 / ResNet-101 vs ViT-Small on VeRi-776}
\end{center}

\end{document}
